% =============================================================================
% 5.4 Discussion of SQ4 — The Crossover Point
% =============================================================================

\section{Discussion of SQ4 --- The Crossover Point}
\label{sec:discuss-sq4}

\textbf{SQ4 asks:} At what concurrency level does the Kubernetes overhead become smaller than the VM's contention-induced degradation --- the crossover point?

\paragraph{Reliability crossover occurs at L3: the primary crossover point.}
The reliability crossover --- the concurrency level at which the VM's success rate first falls below 95\,\% while Kubernetes maintains 100\,\% --- occurred at \textbf{L3 (3 concurrent jobs)}.
This is the primary crossover point for production decision-making: at three concurrent Dagster jobs, the VM already fails one in three runs due to memory exhaustion, while Kubernetes completes all runs successfully within its per-pod limits.

\paragraph{K8s overhead is not flat and grows with load.}
Container startup overhead grows from $\approx$4\,s at L1 to $\approx$15\,s at L10 (mean across levels: 7.9\,s), representing 6--19\,\% of the 63-second job execution time.
Startup overhead increases with concurrency as more containers initialise simultaneously on the single-node cluster, competing for PostgreSQL connections and Python interpreter init.
K8s execution time itself also increases with concurrency: from 63.4\,s (L1) to 81.6\,s (L10), an increase of 29\,\%, as the single-node Kind cluster experiences aggregate memory pressure at higher levels.

\paragraph{The performance crossover is nuanced but the reliability crossover remains decisive.}
In terms of raw wall-clock time per successful job:
K8s is slightly \emph{slower} at L1 (+2.4\,s) and at L7+ (+14.5\,s at L7, +16.6\,s at L10), and \emph{faster} at L2 ($-$7.5\,s), L3 ($-$0.2\,s), and L5 ($-$1.4\,s).
There is no single monotone performance crossover; K8s wall-clock advantage holds at L2--L5 where startup overhead remains low (4--7\,s) and execution times are comparable.
At L7--L10, K8s incurs larger startup overhead (12--15\,s) and higher execution times, making it slower in raw wall-clock time per surviving job; however, the VM delivers fewer than 43\% of those jobs.
The \emph{reliability crossover} at L3 is therefore the dominant and unambiguous signal.

\paragraph{The reliability crossover is the binding constraint.}
At L3, both conditions for migration are met:
\begin{enumerate}
  \item VM success rate falls below 95\,\% (66.7\,\%),
  \item K8s execution time is already competitive ($\approx$69\,s vs 69.9\,s VM survivor mean).
\end{enumerate}
The combined crossover is therefore at \textbf{L3 (3 concurrent jobs)}, and the binding constraint is \emph{reliability}, not performance overhead.

\paragraph{Node-level observations at L10.}
Both environments experience increasing overhead at L10 but from different mechanisms.
The VM fails at the per-container level (OS kills individual containers) from L3 onwards, delivering only 30\% success at L10.
Kubernetes maintained 100\% reliability at L10 in the measured session, with total wall-clock time of 96.5\,s (execution 81.6\,s + startup 14.9\,s).
The K8s reliability advantage over the VM holds across the entire L1--L10 range at this hardware configuration.

\paragraph{The crossover is hardware-configuration dependent.}
The observed crossover at L3 is specific to this configuration:
\begin{itemize}
  \item 4\,vCPU, 4\,GB RAM VM with 2\,GB per-container Docker limit
  \item 4\,vCPU, 8\,GB Kind node with 2\,GiB per-pod memory limit
  \item 60-second two-phase workload with 400\,MB memory allocation per job
  \item 4--15\,s Kubernetes startup overhead (local Kind cluster, pre-pulled images, single node)
\end{itemize}
If the VM had less RAM, the crossover would occur at a lower concurrency level.
In a production K8s cluster (GKE, EKS) with cold image pulls, startup overhead would be higher (10--60\,s), potentially shifting the performance crossover while the reliability crossover remains determined by VM memory exhaustion.

\paragraph{Consistency with the hypothesis.}
The pre-experiment hypothesis (Section~\ref{sec:hypothesis}) predicted that at low concurrency the VM outperforms Kubernetes due to startup overhead, and at high concurrency Kubernetes outperforms the VM.
The data partly confirms this: at L1, the VM has a 2\,s execution time advantage.
However, the hypothesis underestimated how early VM reliability would collapse (L3, not L5--L7) and overestimated K8s startup overhead.
The main finding is that K8s overhead is small enough that the migration decision is primarily a reliability question, not a performance tradeoff.

\paragraph{What the crossover means in practice.}
Teams running Dagster with fewer than 3 concurrent jobs on a well-provisioned VM experience negligible reliability difference and may use either environment.
Teams expecting 3 or more concurrent jobs should migrate to Kubernetes: K8s adds 4--15\,s per-run startup overhead (growing with concurrency) while providing 100\,\% reliability versus the VM's 66.7\,\% or lower success rates at these concurrency levels.
The migration decision threshold is 3 concurrent jobs, and given the negligible overhead, earlier migration is also rational for any team anticipating growth.

